\begin{definition}
Let $\Omega\subset\mathbb{C}$ be open and let $1<p<\infty$.  A mapping
$f\colon\mathbb{C}\to\mathbb{C}$ is said to \emph{satisfy the system (12) on
$\Omega$} if there are weak Wirtinger derivatives $a=\partial f/\partial z$
and $b=\partial f/\partial\bar z$ of $f$ on $\Omega$, in the sense of
\noderef{IsWeakWirtingerOn}, such that
\[
  b(z)
  =\Bigl(\frac1p-\frac12\Bigr)
   \Bigl[\frac{\overline{f(z)}}{f(z)}\,a(z)
     +\frac{f(z)}{\overline{f(z)}}\,\overline{a(z)}\Bigr]
\]
at almost every $z\in\Omega$ with $f(z)\neq0$, ``almost every'' referring to
two-dimensional Lebesgue measure.

This is equation (12) of the paper,
\[
  \frac{\partial f}{\partial\bar z}
  =\Bigl(\frac1p-\frac12\Bigr)
   \Bigl[\frac{\bar f}{f}\frac{\partial f}{\partial z}
        +\frac f{\bar f}\,\overline{\frac{\partial f}{\partial z}}\Bigr],
  \tag{12}
\]
the one complex equation into which Section~2 recasts the $p$-harmonic
equation (1) for the complex gradient $f=\partial u/\partial z$.  The
right-hand side involves the quotients $\bar f/f$ and $f/\bar f$ and is
therefore defined only where $f$ does not vanish, which is why the condition
is imposed off the zero set of $f$; the paper asserts (12) ``almost
everywhere'', and in its situation the zero set of $f$ is discrete, hence of
measure zero, so the two readings agree.

The derivatives are the weak ones.  Equation (12) is introduced at
paper.tex:140--143 for the same $f$, and in the same breath, as the distortion
inequality (11) of \noderef{QuasiregularOn}, and the only hypothesis that $f$
carries there is $f\in W^{1,2}_{\mathrm{loc}}(\Omega)$; so $\partial f/\partial z$
and $\partial f/\partial\bar z$ in (12) are the distributional derivatives,
and satisfying (12) is a property of the class of $f$.  Read through the real
Fr\'echet derivative of a chosen representative, (12) would be satisfied
vacuously by every nowhere differentiable representative of every $L^2$-class,
since both of its sides would vanish identically.  By the next-to-last
paragraph of \noderef{IsWeakWirtingerOn}, whenever $f$ is in addition
differentiable at almost every point of $\Omega$ the pair $(a,b)$ agrees
almost everywhere on $\Omega$ with the pointwise pair $(f_z,f_{\bar z})$ read
through $\mathrm{d}f$, which is how the pointwise form of (12) --- the form
consumed by the elliptic-regularity step and by the hodograph transformation
--- is recovered.  By the last paragraph of \noderef{IsWeakWirtingerOn} the
choice of the witnessing pair is immaterial.
\end{definition}
